\section*{Capítulo \ref{ch:b3:mantle-convection} --- Convecção no manto e forças motrizes das placas}

\begin{solution}{exo:b3:mantle-convection:1}
$\mathrm{Ra} = 3500\times9.9\times\num{3e-5}\times1000\times(\num{6.7e5})^3/
(10^{-6}\times\num{5e20}) = \num{6.25e5}$, cerca de 950 vezes o valor crítico
657.5: a camada convecta vigorosamente.
\end{solution}

\begin{solution}{exo:b3:mantle-convection:2}
$\tau_M = 10^{21}/\num{7.7e10} = \qty{1.30e10}{s}$, cerca de \qty{412}{yr}. O
maré, que dura horas, encontra um manto elástico; o derretimento de um manto de gelo,
ao longo de milhares de anos, e o afundamento de uma laje, ao longo de milhões, encontram um
manto viscoso.
\end{solution}

\begin{solution}{exo:b3:mantle-convection:3}
$\tau = 4\pi\times10^{21}/(3300\times9.8\lambda)$: \qty{3.89e11}{s} =
\qty{12.3}{kyr} para \qty{1000}{km}, e \qty{4.1}{kyr} para \qty{3000}{km}. O
carregamento mais amplo se recupera mais rapidamente.
\end{solution}

\begin{solution}{exo:b3:mantle-convection:4}
$g\rho_m\alpha(T_m - T_0)\kappa t = 9.8\times3300\times\num{3.1e-5}\times1450
\times\num{0.8e-6}\times50\times\num{3.156e13} = \qty{1.835e12}{N.m^{-1}}$; o
colchete é $1 + 2\times3300\times\num{3.1e-5}\times1450/(\pi\times2270) = 1.0416$,
logo $F_{rp} = \qty{1.91e12}{N.m^{-1}}$, 24 vezes menor que a tração da laje.
\end{solution}

\begin{solution}{exo:b3:mantle-convection:5}
(a) $v = 3300\times9.8\times\num{3e-5}\times300\times(\num{e4})^2/10^{21} =
\qty{2.91e-11}{m.s^{-1}}$; o tempo de subida através de \qty{10}{km} é $d/v =
\qty{3.44e14}{s}$, \qty{10.9}{Myr}, o tempo de difusão $d^2/\kappa =
\qty{e14}{s}$, \qty{3.2}{Myr}. A razão entre eles é $(d^2/\kappa)/(d/v) =
vd/\kappa = \rho g\alpha\Delta Td^3/(\kappa\eta) = \mathrm{Ra} = 0.29$. (b) Ele
resfria primeiro: blocos pequenos não convectam. (c) $\mathrm{Ra} \propto d^3$, logo $d =
10\times(657.5/0.291)^{1/3} = \qty{131}{km}$.
\end{solution}

\begin{solution}{exo:b3:mantle-convection:6}
$2\mathrm{Ra}_c = 1315$ e $\mathrm{Ra}/2\mathrm{Ra}_c = 7604$, cuja raiz cúbica
é 19.66. $\delta = 2891/19.66 = \qty{147}{km}$, $\mathrm{Nu} = 19.66/2 =
9.8$, $u = (10^{-6}/\num{2.891e6})\times19.66^2 = \qty{1.34e-10}{m.s^{-1}} =
\qty{0.42}{cm/yr}$, e $q = 9.8\times3.1\times1450/\num{2.891e6} =
\qty{15}{mW.m^{-2}}$.
\end{solution}

\begin{solution}{exo:b3:mantle-convection:7}
$u = 0.09/\num{3.156e7} = \qty{2.85e-9}{m.s^{-1}}$ e $B = 2270\times\num{8e8}
\times\num{2.85e-9} = \qty{5.18e3}{kg.s^{-1}} = \qty{5.2}{Mg/s}$, cerca de
três quintos do valor do Havaí. O fluxo de calor é $1170\times5178/\num{3.1e-5} =
\qty{1.95e11}{W}$, \qty{0.2}{TW}, ou \qty{0.42}{\percent} da perda de
calor terrestre.
\end{solution}

\begin{solution}{exo:b3:mantle-convection:8}
$v = 2\times30\times9.8\times(\num{4e5})^2/(9\times10^{21}) =
\qty{1.05e-8}{m.s^{-1}} = \qty{0.33}{m/yr}$; \qty{2000}{km} levam
$\num{2e6}/\num{1.05e-8} = \qty{1.9e14}{s}$, cerca de \qty{6}{Myr}.
\end{solution}

\begin{solution}{exo:b3:mantle-convection:9}
$w = 300e^{-10/4.4} = \qty{30.9}{m}$ e $\dot w = w/\tau = 30.9/4400 =
\qty{7.0}{mm/yr}$, dois terços da taxa medida perto de Umeå: um
semiespaço uniforme com este $\tau$ está na faixa certa, mas o rebote real é
mais rápido ou sua depressão é mais profunda.
\end{solution}

\begin{solution}{exo:b3:mantle-convection:10}
(a) Com $w = 2\rho_m\alpha\Delta T\sqrt{\kappa t/\pi}/(\rho_m - \rho_w)$,
\[
  \tfrac12(\rho_m - \rho_w)gw^2 = \frac{2g\rho_m^2\alpha^2\Delta T^2\kappa t}
  {\pi(\rho_m - \rho_w)} = g\rho_m\alpha\Delta T\kappa t\times
  \frac{2\rho_m\alpha\Delta T}{\pi(\rho_m - \rho_w)} .
\] (b) Por partes,
$\int_0^\infty x\operatorname{erfc}x\,\mathrm dx = [\tfrac{x^2}2\operatorname{erfc}
x]_0^\infty + \int_0^\infty\tfrac{x^2}2\,\tfrac2{\sqrt\pi}e^{-x^2}\,\mathrm dx =
0 + \tfrac1{\sqrt\pi}\cdot\tfrac{\sqrt\pi}4 = \tfrac14$. (c) Em \qty{100}{Myr},
$\rho_mg\alpha\int z\theta = \qty{3.67e12}{N.m^{-1}}$ e, com $w =
\qty{3705}{m}$, $\tfrac12(\rho_m - \rho_w)gw^2 = \qty{1.53e11}{N.m^{-1}}$: o
termo da água é apenas \qty{4}{\percent} da força.
\end{solution}

\begin{solution}{exo:b3:mantle-convection:11}
$\ln f = 3\ln(\pi^2 + k^2) - 2\ln k$, cuja derivada $6k/(\pi^2 + k^2) - 2/k$
anula-se quando $3k^2 = \pi^2 + k^2$, $k^2 = \pi^2/2$. Nesse caso $f = (3\pi^2/2)^3/
(\pi^2/2) = 27\pi^4/4$. O comprimento de onda $2\pi/k = 2\sqrt2\,d = 2.83\,d$ abrange um
par de células de rotação contrária, cada uma com $\sqrt2\,d = 1.41\,d$ de largura.
\end{solution}

\begin{solution}{exo:b3:mantle-convection:12}
(a) $u = \qty{1.90e-9}{m.s^{-1}}$, $\eta u/h = 10^{20}\times\num{1.90e-9}/
\num{1.5e5} = \qty{1.27}{MPa}$, e ao longo de \qty{6000}{km}, \qty{7.6e12}{N.m^{-1}}.
(b) O dobro do empurrão da dorsal e um sexto da tração da laje: o arrasto pode resistir
ao empurrão da dorsal sozinho, mas não à tração da laje. (c) $\eta uL/h = \num{3.8e12} +
\num{4.6e12}$ dá $u = \num{8.4e12}\times\num{1.5e5}/(\num{6e6}\times10^{20})
= \qty{2.1e-9}{m.s^{-1}} = \qty{6.6}{cm/yr}$, a velocidade das placas oceânicas
rápidas.
\end{solution}

\begin{solution}{pb:b3:mantle-convection:1}
\textbf{1.} $p = 1000\times9.8\times300 = \qty{2.94}{MPa}$.

\textbf{2.} $w = \rho_wh/\rho_m = 300\times1000/3300 = \qty{90.9}{m}$. A crosta
afunda como um bloco; o que escoa de sob a carga, e depois retorna, é o
manto, logo a coluna é equilibrada pelo manto deslocado em profundidade.

\textbf{3.} $\pi\times(\num{1.25e5})^2 = \qty{4.91e10}{m^2}$, e a água
pesa $\num{4.91e10}\times300\times1000 = \qty{1.47e16}{kg}$.

\textbf{4.} A linha de costa formou-se nivelada, quando o fundo estava mais deprimido no
centro; desde então o centro se elevou mais que a borda, logo a linha agora
sobe em direção ao centro, em \qty{77.2}{m} ao longo de \qty{125}{km}:
\qty{0.62}{m/km}.

\textbf{5.} $\lambda/2 = \qty{250}{km}$, logo $\lambda = \qty{500}{km}$ e $k =
2\pi/\num{5e5} = \qty{1.257e-5}{m^{-1}}$.

\textbf{6.} $1/k = \lambda/2\pi = \qty{79.6}{km}$.

\textbf{7.} $w = 90.9 - \text{soerguimento}$: 79.4, 69.3, 52.9, 40.4, 26.9 e
\qty{13.7}{m}.

\textbf{8.} $w = w_0e^{-t/\tau}$ fornece $\ln w = \ln w_0 - t/\tau$, uma reta
de inclinação $-1/\tau$.

\textbf{9.} $\tau = (14 - 1)/\ln(79.4/13.7) = 13/1.7571 = \qty{7.40}{kyr}$.

\textbf{10.} $w(6) = 79.4\,e^{-5/7.40} = \qty{40.4}{m}$, um levantamento de
\qty{50.5}{m}, como na tabela.

\textbf{11.} Restam \qty{13.7}{m}; $\dot w = 13.7/7400 = \qty{1.85}{mm/yr}$.

\textbf{12.} Todos os \qty{13.7}{m}, aproximando-se cada vez mais lentamente.

\textbf{13.} A taxa cai como $e^{-t/\tau}$: $t = 7.40\ln(1.85/0.5) =
7.40\times1.309 = \qty{9.7}{kyr}$, cerca de 9700 anos a mais.

\textbf{14.} $\tau = 7399\times\num{3.156e7} = \qty{2.335e11}{s}$, e $\eta =
\tau\rho_mg\lambda/(4\pi) = \num{2.335e11}\times3300\times9.8\times\num{5e5}/
12.566 = \qty{3.0e20}{Pa.s}$.

\textbf{15.} Logo abaixo da extremidade inferior da faixa fennoscandiana. O fluxo
sob uma carga de \qty{500}{km} de comprimento de onda decai ao longo de
\qty{80}{km}: o rebote amostra o manto mais raso e mais fraco, sob e em parte
dentro do que, na Terra real, é a litosfera.

\textbf{16.} $\tau = 4\pi\times\num{3.0e20}/(3300\times9.8\times\num{3e6}) =
\qty{3.9e10}{s} = \qty{1.2}{kyr}$. Sob uma carga tão larga, o fluxo atinge
\qty{477}{km} de profundidade e mais, em direção a um manto inferior dez vezes
mais viscoso, o que o retarda.

\textbf{17.} Ela sustenta parte de uma carga estreita por flexão: a depressão
em equilíbrio é mais rasa e mais larga que os \qty{90.9}{m} isostáticos, e
parte do rebote é elástica e imediata. Um semiespaço viscoso uniforme ajustado
ao levantamento então fornece uma viscosidade que mistura os dois efeitos.

\textbf{18.} $\tau_M = \num{3.0e20}/\num{7.7e10} = \qty{3.9e9}{s}$, cerca de
\qty{124}{yr}, sessenta vezes menor que $\tau$: ao longo de milhares de anos o
manto responde como um fluido viscoso.

\textbf{19.} $1 - e^{-50/7.40} = 0.9988$: \qty{99.9}{\percent} da depressão de
equilíbrio foi atingida, e a suposição se mantém.

\textbf{20.} $\eta \propto \lambda$: reduzir $\lambda$ à metade reduz $\eta$ à
metade, e $\log_{10}\eta$ diminui de $\log_{10}2 = 0.30$.

\textbf{21.} $\mathrm{Ra} = 3700\times9.9\times\num{3.1e-5}\times1450\times
(\num{6.7e5})^3/(10^{-6}\times\num{3.0e20}) = \num{1.65e6}$.

\textbf{22.} $\num{1.65e6}/657.5 \approx 2500$ vezes o valor crítico.

\textbf{23.} $(\mathrm{Ra}/2\mathrm{Ra}_c)^{1/3} = (1254)^{1/3} = 10.78$: $\delta
= 670/10.78 = \qty{62}{km}$ e $\mathrm{Nu} = 5.4$.

\textbf{24.} $u = (10^{-6}/\num{6.7e5})\times10.78^2 = \qty{1.73e-10}{m.s^{-1}}
= \qty{0.55}{cm/yr}$, de duas a quatorze vezes mais lenta que as placas.

\textbf{25.} $q = 5.4\times3.1\times1450/\num{6.7e5} = \qty{36}{mW.m^{-2}}$. A
camada deixa de lado o calor produzido dentro do manto e o calor do núcleo e do
manto inferior; e a escala descarta suas constantes de ordem um, enquanto as
placas, resfriando-se como semiespaços a partir das dorsais, perdem muito mais
calor em assoalho jovem do que uma camada limite média perderia.

\textbf{26.} $\eta = \qty{3.0e20}{Pa.s}$: $\log_{10}\eta =
\textbf{\num[text-series-to-math=true]{20.48}}$.
\end{solution}
